\documentclass[11pt,a4paper]{article}
\usepackage[utf8]{inputenc}
\usepackage[T1]{fontenc}
\usepackage[spanish,es-minimal,es-noshorthands]{babel}
\AtBeginDocument{\renewcommand{\sin}{\operatorname{sen}}\renewcommand{\arcsin}{\operatorname{arcsen}}}
\usepackage{lmodern}
\usepackage[sfdefault,lining]{FiraSans}
\usepackage[a4paper,top=2.3cm,bottom=2.2cm,left=1.8cm,right=1.8cm,headheight=1.3cm,headsep=0.4cm,footskip=1cm]{geometry}
\usepackage{amsmath,amssymb}
\usepackage{icomma}
\usepackage{graphicx}
\usepackage[table]{xcolor}
\usepackage{tikz}
\usepackage{fontawesome5}
\usepackage{ccicons}
\usepackage{fancyhdr}
\usepackage{lastpage}
\usepackage{tasks}
\usepackage{needspace}
\usepackage[most]{tcolorbox}
\usepackage{array}
\usepackage[hidelinks]{hyperref}
\AtBeginDocument{\hypersetup{pdftitle={Soluciones de la ficha de la unidad 10: Estadística bidimensional. 4º ESO},pdfsubject={Matemáticas A 4º ESO},pdfkeywords={estadística bidimensional, distribuciones bidimensionales, tabla de doble entrada, nube de puntos, centro de gravedad, correlación, covarianza, coeficiente de correlación, recta de regresión, estimaciones}}}

\definecolor{navy}{HTML}{1B2748}
\definecolor{accent}{HTML}{E8603A}
\definecolor{teal}{HTML}{1F9E90}
\definecolor{soft}{HTML}{EEF1F8}
\definecolor{muted}{HTML}{5D6685}

\setlength{\parindent}{0pt}
\setlength{\parskip}{3pt}
\settasks{label=\alph*), label-format=\bfseries\color{navy}, label-width=1.3em, item-indent=1.8em, column-sep=0.6em, after-item-skip=5pt, before-skip=3pt, after-skip=4pt}

% ---------- capçalera i peu ----------
\newcommand{\llicencia}{\href{https://creativecommons.org/licenses/by-nc-sa/4.0/deed.es}{\ccbyncsa}\ \ Licencia Creative Commons Atribución-NoComercial-CompartirIgual 4.0 Internacional (CC BY-NC-SA 4.0)}
\pagestyle{fancy}
\fancyhf{}
\fancyhead[L]{\raisebox{-0.35\height}{\includegraphics[height=0.95cm]{logo-gva.png}}}
\fancyhead[C]{\small\color{muted}Soluciones de la ficha · Unidad 10: Estadística bidimensional · 4\textsuperscript{o} ESO}
\fancyhead[R]{\raisebox{-0.35\height}{\includegraphics[height=0.95cm]{logo-vermellar.png}}}
\fancyfoot[L]{\scriptsize\color{muted}\llicencia}
\fancyfoot[R]{\small\color{muted}\thepage}
\renewcommand{\headrulewidth}{0.4pt}
\renewcommand{\headrule}{\color{muted!40}\hrule height \headrulewidth}
\fancypagestyle{primera}{\fancyhf{}\renewcommand{\headrulewidth}{0pt}
  \fancyfoot[L]{\scriptsize\color{muted}\llicencia}
  \fancyfoot[R]{\small\color{muted}\thepage}}

% ---------- nivells i exercicis ----------
\newcommand{\estrelles}[1]{{\color{accent}\ifcase#1\or\faStar\faStar[regular]\faStar[regular]\or\faStar\faStar\faStar[regular]\or\faStar\faStar\faStar\fi}}
\newcounter{exer}
\newcommand{\ex}[1]{\par\needspace{4\baselineskip}\medskip\stepcounter{exer}\noindent{\color{navy}\bfseries\theexer.}\ \ {\small\estrelles{#1}}\ \ }

\newcommand{\seccio}[2]{\par\needspace{10\baselineskip}\bigskip
  \begin{tcolorbox}[colback=navy,colframe=navy,coltext=white,arc=2mm,boxrule=0pt,left=3mm,right=3mm,top=1.2mm,bottom=1.2mm]
  {\large\bfseries #1\quad #2}
  \end{tcolorbox}\vspace{-2pt}}

\newcolumntype{C}[1]{>{\centering\arraybackslash}p{#1}}
\definecolor{violet}{HTML}{7C6CF0}
\definecolor{gold}{HTML}{E0A21B}
% quadre de recordatori (només a la fitxa d'exercicis)
\newcommand{\recorda}[1]{\par\smallskip\begin{tcolorbox}[colback=accent!7,colframe=accent!55,arc=2mm,boxrule=0.5pt,left=2.5mm,right=2.5mm,top=1mm,bottom=1mm]\small{\color{accent}\faLightbulb[regular]}\ \textbf{Recuerda:}\ #1\end{tcolorbox}}
\newcommand{\res}[1]{\boxed{#1}}\newcommand{\rest}[1]{\fbox{#1}}

\def\fmneed{14}
% =====================================================================
% Capçalera FentMates de les fitxes: un color per a cada curs.
% Els colors són els de la pàgina inicial (LEVELS a assets/js/materials.js);
% el logo ve de _eines/logo/pdf/ (python3 _eines/logo/genera.py).
%
% Ús, al final del preàmbul de la unitat:
%   % =====================================================================
% Capçalera FentMates de les fitxes: un color per a cada curs.
% Els colors són els de la pàgina inicial (LEVELS a assets/js/materials.js);
% el logo ve de _eines/logo/pdf/ (python3 _eines/logo/genera.py).
%
% Ús, al final del preàmbul de la unitat:
%   % =====================================================================
% Capçalera FentMates de les fitxes: un color per a cada curs.
% Els colors són els de la pàgina inicial (LEVELS a assets/js/materials.js);
% el logo ve de _eines/logo/pdf/ (python3 _eines/logo/genera.py).
%
% Ús, al final del preàmbul de la unitat:
%   \input{../comu/fentmates.tex}
%   \fmcurs{1eso}
% i, al començament del document:
%   \capcalera{Matemàtiques}{1r ESO}{Fitxa d'exercicis · Unitat 2}{Divisibilitat}
% Opcional abans de \input: \def\fmneed{14} (espai mínim abans de cada secció).
% =====================================================================
\makeatletter
\def\fm@logodir{../../logo/pdf}
\providecommand{\fmneed}{10}
\def\fm@cursnom{}

% clau del curs, color i color del text damunt del color
\newcommand{\fm@curs}[3]{\@namedef{fm@color@#1}{#2}\@namedef{fm@text@#1}{#3}}
\fm@curs{1eso}{E6B84A}{navy}
\fm@curs{2eso}{3B82F6}{white}
\fm@curs{3eso}{FF7A50}{white}
\fm@curs{4eso-a}{A78BFA}{white}
\fm@curs{4eso-b}{6D28D9}{white}
\fm@curs{1batx-mi}{9B1C1C}{white}
\fm@curs{1batx-ccss}{EF6B6B}{white}
\fm@curs{2batx-mii}{1E6B47}{white}
\fm@curs{2batx-ccss}{2BB3A3}{white}

\newcommand{\fmcurs}[1]{%
  \@ifundefined{fm@color@#1}{\PackageError{fentmates}{Curs desconegut: #1}{}}{}%
  \edef\fm@tmp{\noexpand\definecolor{curs}{HTML}{\@nameuse{fm@color@#1}}%
               \noexpand\colorlet{curstext}{\@nameuse{fm@text@#1}}}\fm@tmp
  \colorlet{cursfosc}{curs!72!black}%
  \def\fm@logo{\fm@logodir/vertical-#1.pdf}%
  \def\fm@marca{\fm@logodir/marca-#1.pdf}}

% ---------- capçalera de la primera pàgina ----------
% #1 matèria (petit), #2 curs, #3 tipus de document i unitat, #4 títol
\newcommand{\capcalera}[4]{%
  \gdef\fm@cursnom{#2}%
  \par\noindent
  \begin{tcolorbox}[colback=curs,colframe=curs,boxrule=0pt,arc=2.2mm,width=3.5cm,height=2.15cm,
      valign=center,left=3mm,right=2mm,top=1mm,bottom=1mm,nobeforeafter,box align=center]
    \color{curstext}{\footnotesize #1}\par\vspace{1pt}{\fontsize{19}{21}\selectfont\bfseries #2\par}%
  \end{tcolorbox}\hspace{5mm}%
  \begin{minipage}[c]{\dimexpr\linewidth-3.5cm-5mm-2.9cm\relax}
    {\color{cursfosc}\small\bfseries\MakeUppercase{#3}}\par\vspace{2pt}
    {\color{navy}\fontsize{20}{23}\selectfont\bfseries #4\par}
  \end{minipage}\hfill
  \begin{minipage}[c]{2.6cm}\raggedleft\includegraphics[width=2.6cm]{\fm@logo}\end{minipage}\par
  \vspace{2pt}}

% ---------- barres de secció del color del curs ----------
\renewcommand{\seccio}[2]{\par\needspace{\fmneed\baselineskip}\bigskip
  \begin{tcolorbox}[colback=curs,colframe=curs,coltext=curstext,arc=2mm,boxrule=0pt,left=3mm,right=3mm,top=1.2mm,bottom=1.2mm]
  {\large\bfseries #1\quad #2}
  \end{tcolorbox}\vspace{-2pt}}

% ---------- estrelles de dificultat del color del curs (a Batxillerat no se'n posen) ----------
\@ifundefined{estrelles}{}{%
  \renewcommand{\estrelles}[1]{{\color{curs}\ifcase#1\or\faStar\faStar[regular]\faStar[regular]\or\faStar\faStar\faStar[regular]\or\faStar\faStar\faStar\fi}}}

% ---------- capçalera i peu de les pàgines ----------
% primera pàgina: logos institucionals al peu
\fancypagestyle{primera}{\fancyhf{}\renewcommand{\headrulewidth}{0pt}
  \fancyfoot[L]{\raisebox{-0.32\height}{\includegraphics[height=0.78cm]{logo-gva.png}}\hspace{5mm}%
    \raisebox{-0.32\height}{\includegraphics[height=0.78cm]{logo-vermellar.png}}}
  \fancyfoot[R]{\scriptsize\color{muted}\href{https://creativecommons.org/licenses/by-nc-sa/4.0/deed.ca}{\ccbyncsa}\ \ CC BY-NC-SA 4.0 · fentmates.org\qquad\small\thepage}}
% la resta: el símbol a l'esquerra i el curs a la dreta, amb el títol al centre com abans
\fancyhead[L]{\raisebox{-0.3\height}{\includegraphics[height=0.62cm]{\fm@marca}}}
\fancyhead[R]{\small\bfseries\color{cursfosc}\fm@cursnom}
\renewcommand{\headrule}{\color{curs!45}\hrule height \headrulewidth}
\makeatother

%   \fmcurs{1eso}
% i, al començament del document:
%   \capcalera{Matemàtiques}{1r ESO}{Fitxa d'exercicis · Unitat 2}{Divisibilitat}
% Opcional abans de \input: \def\fmneed{14} (espai mínim abans de cada secció).
% =====================================================================
\makeatletter
\def\fm@logodir{../../logo/pdf}
\providecommand{\fmneed}{10}
\def\fm@cursnom{}

% clau del curs, color i color del text damunt del color
\newcommand{\fm@curs}[3]{\@namedef{fm@color@#1}{#2}\@namedef{fm@text@#1}{#3}}
\fm@curs{1eso}{E6B84A}{navy}
\fm@curs{2eso}{3B82F6}{white}
\fm@curs{3eso}{FF7A50}{white}
\fm@curs{4eso-a}{A78BFA}{white}
\fm@curs{4eso-b}{6D28D9}{white}
\fm@curs{1batx-mi}{9B1C1C}{white}
\fm@curs{1batx-ccss}{EF6B6B}{white}
\fm@curs{2batx-mii}{1E6B47}{white}
\fm@curs{2batx-ccss}{2BB3A3}{white}

\newcommand{\fmcurs}[1]{%
  \@ifundefined{fm@color@#1}{\PackageError{fentmates}{Curs desconegut: #1}{}}{}%
  \edef\fm@tmp{\noexpand\definecolor{curs}{HTML}{\@nameuse{fm@color@#1}}%
               \noexpand\colorlet{curstext}{\@nameuse{fm@text@#1}}}\fm@tmp
  \colorlet{cursfosc}{curs!72!black}%
  \def\fm@logo{\fm@logodir/vertical-#1.pdf}%
  \def\fm@marca{\fm@logodir/marca-#1.pdf}}

% ---------- capçalera de la primera pàgina ----------
% #1 matèria (petit), #2 curs, #3 tipus de document i unitat, #4 títol
\newcommand{\capcalera}[4]{%
  \gdef\fm@cursnom{#2}%
  \par\noindent
  \begin{tcolorbox}[colback=curs,colframe=curs,boxrule=0pt,arc=2.2mm,width=3.5cm,height=2.15cm,
      valign=center,left=3mm,right=2mm,top=1mm,bottom=1mm,nobeforeafter,box align=center]
    \color{curstext}{\footnotesize #1}\par\vspace{1pt}{\fontsize{19}{21}\selectfont\bfseries #2\par}%
  \end{tcolorbox}\hspace{5mm}%
  \begin{minipage}[c]{\dimexpr\linewidth-3.5cm-5mm-2.9cm\relax}
    {\color{cursfosc}\small\bfseries\MakeUppercase{#3}}\par\vspace{2pt}
    {\color{navy}\fontsize{20}{23}\selectfont\bfseries #4\par}
  \end{minipage}\hfill
  \begin{minipage}[c]{2.6cm}\raggedleft\includegraphics[width=2.6cm]{\fm@logo}\end{minipage}\par
  \vspace{2pt}}

% ---------- barres de secció del color del curs ----------
\renewcommand{\seccio}[2]{\par\needspace{\fmneed\baselineskip}\bigskip
  \begin{tcolorbox}[colback=curs,colframe=curs,coltext=curstext,arc=2mm,boxrule=0pt,left=3mm,right=3mm,top=1.2mm,bottom=1.2mm]
  {\large\bfseries #1\quad #2}
  \end{tcolorbox}\vspace{-2pt}}

% ---------- estrelles de dificultat del color del curs (a Batxillerat no se'n posen) ----------
\@ifundefined{estrelles}{}{%
  \renewcommand{\estrelles}[1]{{\color{curs}\ifcase#1\or\faStar\faStar[regular]\faStar[regular]\or\faStar\faStar\faStar[regular]\or\faStar\faStar\faStar\fi}}}

% ---------- capçalera i peu de les pàgines ----------
% primera pàgina: logos institucionals al peu
\fancypagestyle{primera}{\fancyhf{}\renewcommand{\headrulewidth}{0pt}
  \fancyfoot[L]{\raisebox{-0.32\height}{\includegraphics[height=0.78cm]{logo-gva.png}}\hspace{5mm}%
    \raisebox{-0.32\height}{\includegraphics[height=0.78cm]{logo-vermellar.png}}}
  \fancyfoot[R]{\scriptsize\color{muted}\href{https://creativecommons.org/licenses/by-nc-sa/4.0/deed.ca}{\ccbyncsa}\ \ CC BY-NC-SA 4.0 · fentmates.org\qquad\small\thepage}}
% la resta: el símbol a l'esquerra i el curs a la dreta, amb el títol al centre com abans
\fancyhead[L]{\raisebox{-0.3\height}{\includegraphics[height=0.62cm]{\fm@marca}}}
\fancyhead[R]{\small\bfseries\color{cursfosc}\fm@cursnom}
\renewcommand{\headrule}{\color{curs!45}\hrule height \headrulewidth}
\makeatother

%   \fmcurs{1eso}
% i, al començament del document:
%   \capcalera{Matemàtiques}{1r ESO}{Fitxa d'exercicis · Unitat 2}{Divisibilitat}
% Opcional abans de \input: \def\fmneed{14} (espai mínim abans de cada secció).
% =====================================================================
\makeatletter
\def\fm@logodir{../../logo/pdf}
\providecommand{\fmneed}{10}
\def\fm@cursnom{}

% clau del curs, color i color del text damunt del color
\newcommand{\fm@curs}[3]{\@namedef{fm@color@#1}{#2}\@namedef{fm@text@#1}{#3}}
\fm@curs{1eso}{E6B84A}{navy}
\fm@curs{2eso}{3B82F6}{white}
\fm@curs{3eso}{FF7A50}{white}
\fm@curs{4eso-a}{A78BFA}{white}
\fm@curs{4eso-b}{6D28D9}{white}
\fm@curs{1batx-mi}{9B1C1C}{white}
\fm@curs{1batx-ccss}{EF6B6B}{white}
\fm@curs{2batx-mii}{1E6B47}{white}
\fm@curs{2batx-ccss}{2BB3A3}{white}

\newcommand{\fmcurs}[1]{%
  \@ifundefined{fm@color@#1}{\PackageError{fentmates}{Curs desconegut: #1}{}}{}%
  \edef\fm@tmp{\noexpand\definecolor{curs}{HTML}{\@nameuse{fm@color@#1}}%
               \noexpand\colorlet{curstext}{\@nameuse{fm@text@#1}}}\fm@tmp
  \colorlet{cursfosc}{curs!72!black}%
  \def\fm@logo{\fm@logodir/vertical-#1.pdf}%
  \def\fm@marca{\fm@logodir/marca-#1.pdf}}

% ---------- capçalera de la primera pàgina ----------
% #1 matèria (petit), #2 curs, #3 tipus de document i unitat, #4 títol
\newcommand{\capcalera}[4]{%
  \gdef\fm@cursnom{#2}%
  \par\noindent
  \begin{tcolorbox}[colback=curs,colframe=curs,boxrule=0pt,arc=2.2mm,width=3.5cm,height=2.15cm,
      valign=center,left=3mm,right=2mm,top=1mm,bottom=1mm,nobeforeafter,box align=center]
    \color{curstext}{\footnotesize #1}\par\vspace{1pt}{\fontsize{19}{21}\selectfont\bfseries #2\par}%
  \end{tcolorbox}\hspace{5mm}%
  \begin{minipage}[c]{\dimexpr\linewidth-3.5cm-5mm-2.9cm\relax}
    {\color{cursfosc}\small\bfseries\MakeUppercase{#3}}\par\vspace{2pt}
    {\color{navy}\fontsize{20}{23}\selectfont\bfseries #4\par}
  \end{minipage}\hfill
  \begin{minipage}[c]{2.6cm}\raggedleft\includegraphics[width=2.6cm]{\fm@logo}\end{minipage}\par
  \vspace{2pt}}

% ---------- barres de secció del color del curs ----------
\renewcommand{\seccio}[2]{\par\needspace{\fmneed\baselineskip}\bigskip
  \begin{tcolorbox}[colback=curs,colframe=curs,coltext=curstext,arc=2mm,boxrule=0pt,left=3mm,right=3mm,top=1.2mm,bottom=1.2mm]
  {\large\bfseries #1\quad #2}
  \end{tcolorbox}\vspace{-2pt}}

% ---------- estrelles de dificultat del color del curs (a Batxillerat no se'n posen) ----------
\@ifundefined{estrelles}{}{%
  \renewcommand{\estrelles}[1]{{\color{curs}\ifcase#1\or\faStar\faStar[regular]\faStar[regular]\or\faStar\faStar\faStar[regular]\or\faStar\faStar\faStar\fi}}}

% ---------- capçalera i peu de les pàgines ----------
% primera pàgina: logos institucionals al peu
\fancypagestyle{primera}{\fancyhf{}\renewcommand{\headrulewidth}{0pt}
  \fancyfoot[L]{\raisebox{-0.32\height}{\includegraphics[height=0.78cm]{logo-gva.png}}\hspace{5mm}%
    \raisebox{-0.32\height}{\includegraphics[height=0.78cm]{logo-vermellar.png}}}
  \fancyfoot[R]{\scriptsize\color{muted}\href{https://creativecommons.org/licenses/by-nc-sa/4.0/deed.ca}{\ccbyncsa}\ \ CC BY-NC-SA 4.0 · fentmates.org\qquad\small\thepage}}
% la resta: el símbol a l'esquerra i el curs a la dreta, amb el títol al centre com abans
\fancyhead[L]{\raisebox{-0.3\height}{\includegraphics[height=0.62cm]{\fm@marca}}}
\fancyhead[R]{\small\bfseries\color{cursfosc}\fm@cursnom}
\renewcommand{\headrule}{\color{curs!45}\hrule height \headrulewidth}
\makeatother

\fmcurs{4eso-a}

\begin{document}
\thispagestyle{primera}

\capcalera{Matemáticas A}{4\textsuperscript{o} ESO}{Soluciones de la ficha · Unidad 10}{Estadística bidimensional}
\vspace{8pt}

\begin{tcolorbox}[colback=soft,colframe=soft,arc=2mm,left=3mm,right=3mm,top=1.5mm,bottom=1.5mm]
\small Intenta resolver cada ejercicio antes de mirar su solución. Si no te sale el mismo resultado, repasa los pasos o consulta los apuntes interactivos de la unidad: encontrarás herramientas para hacer tablas de doble entrada, dibujar nubes de puntos, calcular el coeficiente de correlación y la recta de regresión paso a paso y hacer estimaciones. Los resultados aproximados están redondeados a las centésimas; haz los cálculos con todos los decimales y redondea solo al final.
\end{tcolorbox}

% =====================================================
\seccio{1}{Distribuciones bidimensionales}

\ex{1} \begin{tasks}(2)
  \task \rest{Funcional}: $A=\pi r^2$.
  \task \rest{Estadística}.
  \task \rest{Ninguna relación}.
  \task \rest{Funcional}: $y=2+1{,}1x$.
  \task \rest{Estadística}.
  \task \rest{Funcional}: $y=2026-x$, por ejemplo.
\end{tasks}

\ex{1} \begin{tasks}(1)
  \task $(1,\,10),\ \allowbreak (2,\,12),\ \allowbreak (2{,}5,\,15),\ \allowbreak (3,\,14),\ \allowbreak (4,\,20),\ \allowbreak (5,\,22),\ \allowbreak (6,\,25),\ \allowbreak (8,\,30)$; \ $N=\res{8}$.
  \task A más de 3 km: $\res{4}$ alumnos. Más de 15 minutos: $\res{4}$ alumnos.
  \task \rest{Estadística}: los que viven más lejos suelen tardar más, pero no siempre (quien vive a 2,5 km tarda más que quien vive a 3 km).
\end{tasks}

\ex{1} \begin{tasks}(2)
  \task $\res{3}$
  \task $\res{10}$
  \task $3+6+2=\res{11}$
  \task $\frac{8}{30}\approx\res{26{,}7\,\%}$
\end{tasks}

\ex{2} \par\smallskip\begin{center}{\renewcommand{\arraystretch}{1.25}\begin{tabular}{|c|c|c|c|c|}\hline $X\backslash Y$ & $0$ & $1$ & $2$ & $f_x$\\\hline $1$ & $2$ & $3$ & $0$ & $5$\\\hline $2$ & $1$ & $4$ & $2$ & $7$\\\hline $3$ & $0$ & $1$ & $3$ & $4$\\\hline $f_y$ & $3$ & $8$ & $5$ & $16$\\\hline\end{tabular}}\end{center}

\ex{2} $\bar x=\frac{0\cdot8+1\cdot14+2\cdot8}{30}=\frac{30}{30}=\res{1}$ hermano; \ $\bar y=\frac{0\cdot10+1\cdot14+2\cdot6}{30}=\frac{26}{30}\approx\res{0{,}87}$ mascotas.

\ex{3} Fila 1: $7-4-1=\res{2}$. Fila 2: $10-5-2=\res{3}$. Fila 3: $8-1-3=\res{4}$. Total de la columna $y=0$: $4+3+1=\res{8}$. Total: $7+10+8=\res{25}$ (o bien $8+10+7$).\par\smallskip\begin{center}{\renewcommand{\arraystretch}{1.25}\begin{tabular}{|c|c|c|c|c|}\hline $X\backslash Y$ & $0$ & $1$ & $2$ & $f_x$\\\hline $1$ & $4$ & $2$ & $1$ & $7$\\\hline $2$ & $3$ & $5$ & $2$ & $10$\\\hline $3$ & $1$ & $3$ & $4$ & $8$\\\hline $f_y$ & $8$ & $10$ & $7$ & $25$\\\hline\end{tabular}}\end{center}


% =====================================================
\seccio{2}{Nube de puntos}

\ex{1} $\bar x=\frac{21}{6}=3{,}5$, \ $\bar y=\frac{27}{6}=4{,}5$: \ centro de gravedad $\res{(3{,}5,\ 4{,}5)}$.\par\smallskip\begin{center}\begin{tikzpicture}[baseline=(current bounding box.center),line cap=round]
\draw[muted!25,line width=0.3pt] (0.000,0)--(0.000,3.600); \node[below,font=\tiny,text=muted,inner sep=1.5pt] at (0.000,0) {$0$};
\draw[muted!25,line width=0.3pt] (0.714,0)--(0.714,3.600); \node[below,font=\tiny,text=muted,inner sep=1.5pt] at (0.714,0) {$1$};
\draw[muted!25,line width=0.3pt] (1.429,0)--(1.429,3.600); \node[below,font=\tiny,text=muted,inner sep=1.5pt] at (1.429,0) {$2$};
\draw[muted!25,line width=0.3pt] (2.143,0)--(2.143,3.600); \node[below,font=\tiny,text=muted,inner sep=1.5pt] at (2.143,0) {$3$};
\draw[muted!25,line width=0.3pt] (2.857,0)--(2.857,3.600); \node[below,font=\tiny,text=muted,inner sep=1.5pt] at (2.857,0) {$4$};
\draw[muted!25,line width=0.3pt] (3.571,0)--(3.571,3.600); \node[below,font=\tiny,text=muted,inner sep=1.5pt] at (3.571,0) {$5$};
\draw[muted!25,line width=0.3pt] (4.286,0)--(4.286,3.600); \node[below,font=\tiny,text=muted,inner sep=1.5pt] at (4.286,0) {$6$};
\draw[muted!25,line width=0.3pt] (5.000,0)--(5.000,3.600); \node[below,font=\tiny,text=muted,inner sep=1.5pt] at (5.000,0) {$7$};
\draw[muted!25,line width=0.3pt] (0,0.000)--(5.000,0.000); \node[left,font=\tiny,text=muted,inner sep=1.5pt] at (0,0.000) {$0$};
\draw[muted!25,line width=0.3pt] (0,0.450)--(5.000,0.450); \node[left,font=\tiny,text=muted,inner sep=1.5pt] at (0,0.450) {$1$};
\draw[muted!25,line width=0.3pt] (0,0.900)--(5.000,0.900); \node[left,font=\tiny,text=muted,inner sep=1.5pt] at (0,0.900) {$2$};
\draw[muted!25,line width=0.3pt] (0,1.350)--(5.000,1.350); \node[left,font=\tiny,text=muted,inner sep=1.5pt] at (0,1.350) {$3$};
\draw[muted!25,line width=0.3pt] (0,1.800)--(5.000,1.800); \node[left,font=\tiny,text=muted,inner sep=1.5pt] at (0,1.800) {$4$};
\draw[muted!25,line width=0.3pt] (0,2.250)--(5.000,2.250); \node[left,font=\tiny,text=muted,inner sep=1.5pt] at (0,2.250) {$5$};
\draw[muted!25,line width=0.3pt] (0,2.700)--(5.000,2.700); \node[left,font=\tiny,text=muted,inner sep=1.5pt] at (0,2.700) {$6$};
\draw[muted!25,line width=0.3pt] (0,3.150)--(5.000,3.150); \node[left,font=\tiny,text=muted,inner sep=1.5pt] at (0,3.150) {$7$};
\draw[muted!25,line width=0.3pt] (0,3.600)--(5.000,3.600); \node[left,font=\tiny,text=muted,inner sep=1.5pt] at (0,3.600) {$8$};
\draw[-latex,muted,line width=0.6pt] (0,0)--(5.300,0) node[right,font=\scriptsize,text=navy]{X};
\draw[-latex,muted,line width=0.6pt] (0,0)--(0,3.900) node[above,font=\scriptsize,text=navy]{Y};
\fill[teal] (0.714,0.900) circle (0.085);
\fill[teal] (1.429,1.800) circle (0.085);
\fill[teal] (2.143,1.350) circle (0.085);
\fill[teal] (2.857,2.700) circle (0.085);
\fill[teal] (3.571,2.250) circle (0.085);
\fill[teal] (4.286,3.150) circle (0.085);
\draw[navy,line width=1pt,fill=white] (2.500,2.025) circle (0.09);
\end{tikzpicture}\end{center}

\ex{1} \par\smallskip\begin{center}{\renewcommand{\arraystretch}{1.3}\begin{tabular}{|l|C{0.8cm}|C{0.8cm}|C{0.8cm}|C{0.8cm}|C{0.8cm}|}\hline $x_i$ & $1$ & $2$ & $4$ & $5$ & $8$\\\hline $y_i$ & $8$ & $6$ & $5$ & $3$ & $3$\\\hline\end{tabular}}\end{center}$\bar x=\frac{20}{5}=4$, \ $\bar y=\frac{25}{5}=5$: \ $\res{(4,\ 5)}$.

\ex{2} $\bar x=\frac{975}{6}=162{,}5$ cm, \ $\bar y=\frac{344}{6}\approx57{,}33$ kg: \ $\res{(162{,}5;\ 57{,}33)}$.\par\smallskip\begin{center}\begin{tikzpicture}[baseline=(current bounding box.center),line cap=round]
\draw[muted!25,line width=0.3pt] (0.000,0)--(0.000,3.600); \node[below,font=\tiny,text=muted,inner sep=1.5pt] at (0.000,0) {$145$};
\draw[muted!25,line width=0.3pt] (0.714,0)--(0.714,3.600); \node[below,font=\tiny,text=muted,inner sep=1.5pt] at (0.714,0) {$150$};
\draw[muted!25,line width=0.3pt] (1.429,0)--(1.429,3.600); \node[below,font=\tiny,text=muted,inner sep=1.5pt] at (1.429,0) {$155$};
\draw[muted!25,line width=0.3pt] (2.143,0)--(2.143,3.600); \node[below,font=\tiny,text=muted,inner sep=1.5pt] at (2.143,0) {$160$};
\draw[muted!25,line width=0.3pt] (2.857,0)--(2.857,3.600); \node[below,font=\tiny,text=muted,inner sep=1.5pt] at (2.857,0) {$165$};
\draw[muted!25,line width=0.3pt] (3.571,0)--(3.571,3.600); \node[below,font=\tiny,text=muted,inner sep=1.5pt] at (3.571,0) {$170$};
\draw[muted!25,line width=0.3pt] (4.286,0)--(4.286,3.600); \node[below,font=\tiny,text=muted,inner sep=1.5pt] at (4.286,0) {$175$};
\draw[muted!25,line width=0.3pt] (5.000,0)--(5.000,3.600); \node[below,font=\tiny,text=muted,inner sep=1.5pt] at (5.000,0) {$180$};
\draw[muted!25,line width=0.3pt] (0,0.000)--(5.000,0.000); \node[left,font=\tiny,text=muted,inner sep=1.5pt] at (0,0.000) {$40$};
\draw[muted!25,line width=0.3pt] (0,0.514)--(5.000,0.514); \node[left,font=\tiny,text=muted,inner sep=1.5pt] at (0,0.514) {$45$};
\draw[muted!25,line width=0.3pt] (0,1.029)--(5.000,1.029); \node[left,font=\tiny,text=muted,inner sep=1.5pt] at (0,1.029) {$50$};
\draw[muted!25,line width=0.3pt] (0,1.543)--(5.000,1.543); \node[left,font=\tiny,text=muted,inner sep=1.5pt] at (0,1.543) {$55$};
\draw[muted!25,line width=0.3pt] (0,2.057)--(5.000,2.057); \node[left,font=\tiny,text=muted,inner sep=1.5pt] at (0,2.057) {$60$};
\draw[muted!25,line width=0.3pt] (0,2.571)--(5.000,2.571); \node[left,font=\tiny,text=muted,inner sep=1.5pt] at (0,2.571) {$65$};
\draw[muted!25,line width=0.3pt] (0,3.086)--(5.000,3.086); \node[left,font=\tiny,text=muted,inner sep=1.5pt] at (0,3.086) {$70$};
\draw[muted!25,line width=0.3pt] (0,3.600)--(5.000,3.600); \node[left,font=\tiny,text=muted,inner sep=1.5pt] at (0,3.600) {$75$};
\draw[-latex,muted,line width=0.6pt] (0,0)--(5.300,0) node[right,font=\scriptsize,text=navy]{cm};
\draw[-latex,muted,line width=0.6pt] (0,0)--(0,3.900) node[above,font=\scriptsize,text=navy]{kg};
\fill[teal] (0.714,0.514) circle (0.085);
\fill[teal] (1.429,1.234) circle (0.085);
\fill[teal] (2.143,1.440) circle (0.085);
\fill[teal] (2.857,2.057) circle (0.085);
\fill[teal] (3.571,2.366) circle (0.085);
\fill[teal] (4.286,3.086) circle (0.085);
\draw[navy,line width=1pt,fill=white] (2.500,1.783) circle (0.09);
\end{tikzpicture}\end{center}

\ex{2} \par\smallskip\begin{center}\begin{tikzpicture}[baseline=(current bounding box.center),line cap=round]
\draw[muted!25,line width=0.3pt] (0.000,0)--(0.000,3.000); \node[below,font=\tiny,text=muted,inner sep=1.5pt] at (0.000,0) {$0$};
\draw[muted!25,line width=0.3pt] (1.067,0)--(1.067,3.000); \node[below,font=\tiny,text=muted,inner sep=1.5pt] at (1.067,0) {$1$};
\draw[muted!25,line width=0.3pt] (2.133,0)--(2.133,3.000); \node[below,font=\tiny,text=muted,inner sep=1.5pt] at (2.133,0) {$2$};
\draw[muted!25,line width=0.3pt] (3.200,0)--(3.200,3.000); \node[below,font=\tiny,text=muted,inner sep=1.5pt] at (3.200,0) {$3$};
\draw[muted!25,line width=0.3pt] (0,0.000)--(3.200,0.000); \node[left,font=\tiny,text=muted,inner sep=1.5pt] at (0,0.000) {$0$};
\draw[muted!25,line width=0.3pt] (0,1.000)--(3.200,1.000); \node[left,font=\tiny,text=muted,inner sep=1.5pt] at (0,1.000) {$1$};
\draw[muted!25,line width=0.3pt] (0,2.000)--(3.200,2.000); \node[left,font=\tiny,text=muted,inner sep=1.5pt] at (0,2.000) {$2$};
\draw[muted!25,line width=0.3pt] (0,3.000)--(3.200,3.000); \node[left,font=\tiny,text=muted,inner sep=1.5pt] at (0,3.000) {$3$};
\draw[-latex,muted,line width=0.6pt] (0,0)--(3.500,0) node[right,font=\scriptsize,text=navy]{X};
\draw[-latex,muted,line width=0.6pt] (0,0)--(0,3.300) node[above,font=\scriptsize,text=navy]{Y};
\fill[teal] (0.000,0.000) circle (0.085);
\node[above right,font=\tiny,inner sep=1pt] at (0.000,0.000) {$3$};
\fill[teal] (0.000,1.000) circle (0.085);
\node[above right,font=\tiny,inner sep=1pt] at (0.000,1.000) {$4$};
\fill[teal] (0.000,2.000) circle (0.085);
\node[above right,font=\tiny,inner sep=1pt] at (0.000,2.000) {$1$};
\fill[teal] (1.067,0.000) circle (0.085);
\node[above right,font=\tiny,inner sep=1pt] at (1.067,0.000) {$5$};
\fill[teal] (1.067,1.000) circle (0.085);
\node[above right,font=\tiny,inner sep=1pt] at (1.067,1.000) {$6$};
\fill[teal] (1.067,2.000) circle (0.085);
\node[above right,font=\tiny,inner sep=1pt] at (1.067,2.000) {$3$};
\fill[teal] (2.133,0.000) circle (0.085);
\node[above right,font=\tiny,inner sep=1pt] at (2.133,0.000) {$2$};
\fill[teal] (2.133,1.000) circle (0.085);
\node[above right,font=\tiny,inner sep=1pt] at (2.133,1.000) {$4$};
\fill[teal] (2.133,2.000) circle (0.085);
\node[above right,font=\tiny,inner sep=1pt] at (2.133,2.000) {$2$};
\end{tikzpicture}\end{center}Con las medias del ejercicio 5: $\res{(1;\ 0{,}87)}$.

\ex{2} Las $x$ tienen que sumar $5\cdot4{,}2=21$ y las $y$, $5\cdot5=25$: \ $x=21-16=5$, \ $y=25-20=5$: \ $\res{(5,\ 5)}$.

\ex{3} Ahora es $(\frac{12}{3},\frac{15}{3})=(4,\ 5)$. Si añadimos el mismo punto $\res{(4,\ 5)}$, las medias no cambian. Para $(5,\ 6)$ con 4 puntos, las $x$ tienen que sumar $20$ y las $y$, $24$: hay que añadir $\res{(8,\ 9)}$.


% =====================================================
\seccio{3}{Correlación}

\ex{1} \begin{tasks}(3)
  \task \rest{Positiva y fuerte}
  \task \rest{Negativa y débil}
  \task \rest{Nula}
  \task \rest{Positiva y débil}
  \task \rest{Negativa y fuerte}
  \task \rest{No lineal}: los puntos siguen una curva (parábola).
\end{tasks}

\ex{1} \begin{tasks}(2)
  \task \rest{Negativa}
  \task \rest{Negativa} (menos tiempo)
  \task \rest{Positiva}
  \task \rest{Positiva}
  \task \rest{Nula}
  \task \rest{Negativa}
\end{tasks}

\ex{1} \begin{tasks}(4)
  \task A: $\res{r=-0{,}9}$
  \task B: $\res{r=0{,}4}$
  \task C: $\res{r=0{,}9}$
  \task D: $\res{r=0}$
\end{tasks}

\ex{2} \par\smallskip\begin{center}\begin{tikzpicture}[baseline=(current bounding box.center),line cap=round]
\draw[muted!25,line width=0.3pt] (0.000,0)--(0.000,3.600); \node[below,font=\tiny,text=muted,inner sep=1.5pt] at (0.000,0) {$0$};
\draw[muted!25,line width=0.3pt] (0.833,0)--(0.833,3.600); \node[below,font=\tiny,text=muted,inner sep=1.5pt] at (0.833,0) {$1$};
\draw[muted!25,line width=0.3pt] (1.667,0)--(1.667,3.600); \node[below,font=\tiny,text=muted,inner sep=1.5pt] at (1.667,0) {$2$};
\draw[muted!25,line width=0.3pt] (2.500,0)--(2.500,3.600); \node[below,font=\tiny,text=muted,inner sep=1.5pt] at (2.500,0) {$3$};
\draw[muted!25,line width=0.3pt] (3.333,0)--(3.333,3.600); \node[below,font=\tiny,text=muted,inner sep=1.5pt] at (3.333,0) {$4$};
\draw[muted!25,line width=0.3pt] (4.167,0)--(4.167,3.600); \node[below,font=\tiny,text=muted,inner sep=1.5pt] at (4.167,0) {$5$};
\draw[muted!25,line width=0.3pt] (5.000,0)--(5.000,3.600); \node[below,font=\tiny,text=muted,inner sep=1.5pt] at (5.000,0) {$6$};
\draw[muted!25,line width=0.3pt] (0,0.000)--(5.000,0.000); \node[left,font=\tiny,text=muted,inner sep=1.5pt] at (0,0.000) {$0$};
\draw[muted!25,line width=0.3pt] (0,0.600)--(5.000,0.600); \node[left,font=\tiny,text=muted,inner sep=1.5pt] at (0,0.600) {$10$};
\draw[muted!25,line width=0.3pt] (0,1.200)--(5.000,1.200); \node[left,font=\tiny,text=muted,inner sep=1.5pt] at (0,1.200) {$20$};
\draw[muted!25,line width=0.3pt] (0,1.800)--(5.000,1.800); \node[left,font=\tiny,text=muted,inner sep=1.5pt] at (0,1.800) {$30$};
\draw[muted!25,line width=0.3pt] (0,2.400)--(5.000,2.400); \node[left,font=\tiny,text=muted,inner sep=1.5pt] at (0,2.400) {$40$};
\draw[muted!25,line width=0.3pt] (0,3.000)--(5.000,3.000); \node[left,font=\tiny,text=muted,inner sep=1.5pt] at (0,3.000) {$50$};
\draw[muted!25,line width=0.3pt] (0,3.600)--(5.000,3.600); \node[left,font=\tiny,text=muted,inner sep=1.5pt] at (0,3.600) {$60$};
\draw[-latex,muted,line width=0.6pt] (0,0)--(5.300,0) node[right,font=\scriptsize,text=navy]{€};
\draw[-latex,muted,line width=0.6pt] (0,0)--(0,3.900) node[above,font=\scriptsize,text=navy]{unidades};
\fill[teal] (0.833,3.000) circle (0.085);
\fill[teal] (1.667,2.640) circle (0.085);
\fill[teal] (2.500,2.400) circle (0.085);
\fill[teal] (3.333,1.920) circle (0.085);
\fill[teal] (4.167,1.440) circle (0.085);
\end{tikzpicture}\end{center}Los puntos bajan y están casi en una recta: correlación \rest{negativa y fuerte}.

\ex{2} \rest{No}. Hay correlación positiva, pero no es causa-efecto: las dos variables dependen de una tercera, el \textbf{tamaño del incendio}. Cuanto más grande es, más bomberos van y más daños hay.

\ex{3} Por ejemplo, negativa y fuerte: $(1,\,9),\ (2,\,8),\ (3,\,6),\ (4,\,5),\ (5,\,2)$ (los puntos bajan casi en línea recta). Nula: $(1,\,3),\ (2,\,7),\ (3,\,1),\ (4,\,6),\ (5,\,3)$ (no suben ni bajan). Hay muchas otras respuestas.\par\smallskip\begin{center}\begin{tikzpicture}[baseline=(current bounding box.center),line cap=round]
\draw[muted!25,line width=0.3pt] (0.000,0)--(0.000,2.600); \node[below,font=\tiny,text=muted,inner sep=1.5pt] at (0.000,0) {$0$};
\draw[muted!25,line width=0.3pt] (0.533,0)--(0.533,2.600); \node[below,font=\tiny,text=muted,inner sep=1.5pt] at (0.533,0) {$1$};
\draw[muted!25,line width=0.3pt] (1.067,0)--(1.067,2.600); \node[below,font=\tiny,text=muted,inner sep=1.5pt] at (1.067,0) {$2$};
\draw[muted!25,line width=0.3pt] (1.600,0)--(1.600,2.600); \node[below,font=\tiny,text=muted,inner sep=1.5pt] at (1.600,0) {$3$};
\draw[muted!25,line width=0.3pt] (2.133,0)--(2.133,2.600); \node[below,font=\tiny,text=muted,inner sep=1.5pt] at (2.133,0) {$4$};
\draw[muted!25,line width=0.3pt] (2.667,0)--(2.667,2.600); \node[below,font=\tiny,text=muted,inner sep=1.5pt] at (2.667,0) {$5$};
\draw[muted!25,line width=0.3pt] (3.200,0)--(3.200,2.600); \node[below,font=\tiny,text=muted,inner sep=1.5pt] at (3.200,0) {$6$};
\draw[muted!25,line width=0.3pt] (0,0.000)--(3.200,0.000); \node[left,font=\tiny,text=muted,inner sep=1.5pt] at (0,0.000) {$0$};
\draw[muted!25,line width=0.3pt] (0,0.520)--(3.200,0.520); \node[left,font=\tiny,text=muted,inner sep=1.5pt] at (0,0.520) {$2$};
\draw[muted!25,line width=0.3pt] (0,1.040)--(3.200,1.040); \node[left,font=\tiny,text=muted,inner sep=1.5pt] at (0,1.040) {$4$};
\draw[muted!25,line width=0.3pt] (0,1.560)--(3.200,1.560); \node[left,font=\tiny,text=muted,inner sep=1.5pt] at (0,1.560) {$6$};
\draw[muted!25,line width=0.3pt] (0,2.080)--(3.200,2.080); \node[left,font=\tiny,text=muted,inner sep=1.5pt] at (0,2.080) {$8$};
\draw[muted!25,line width=0.3pt] (0,2.600)--(3.200,2.600); \node[left,font=\tiny,text=muted,inner sep=1.5pt] at (0,2.600) {$10$};
\draw[-latex,muted,line width=0.6pt] (0,0)--(3.500,0) node[right,font=\scriptsize,text=navy]{X};
\draw[-latex,muted,line width=0.6pt] (0,0)--(0,2.900) node[above,font=\scriptsize,text=navy]{Y};
\fill[teal] (0.533,2.340) circle (0.085);
\fill[teal] (1.067,2.080) circle (0.085);
\fill[teal] (1.600,1.560) circle (0.085);
\fill[teal] (2.133,1.300) circle (0.085);
\fill[teal] (2.667,0.520) circle (0.085);
\end{tikzpicture}\hspace{0.6cm}\begin{tikzpicture}[baseline=(current bounding box.center),line cap=round]
\draw[muted!25,line width=0.3pt] (0.000,0)--(0.000,2.600); \node[below,font=\tiny,text=muted,inner sep=1.5pt] at (0.000,0) {$0$};
\draw[muted!25,line width=0.3pt] (0.533,0)--(0.533,2.600); \node[below,font=\tiny,text=muted,inner sep=1.5pt] at (0.533,0) {$1$};
\draw[muted!25,line width=0.3pt] (1.067,0)--(1.067,2.600); \node[below,font=\tiny,text=muted,inner sep=1.5pt] at (1.067,0) {$2$};
\draw[muted!25,line width=0.3pt] (1.600,0)--(1.600,2.600); \node[below,font=\tiny,text=muted,inner sep=1.5pt] at (1.600,0) {$3$};
\draw[muted!25,line width=0.3pt] (2.133,0)--(2.133,2.600); \node[below,font=\tiny,text=muted,inner sep=1.5pt] at (2.133,0) {$4$};
\draw[muted!25,line width=0.3pt] (2.667,0)--(2.667,2.600); \node[below,font=\tiny,text=muted,inner sep=1.5pt] at (2.667,0) {$5$};
\draw[muted!25,line width=0.3pt] (3.200,0)--(3.200,2.600); \node[below,font=\tiny,text=muted,inner sep=1.5pt] at (3.200,0) {$6$};
\draw[muted!25,line width=0.3pt] (0,0.000)--(3.200,0.000); \node[left,font=\tiny,text=muted,inner sep=1.5pt] at (0,0.000) {$0$};
\draw[muted!25,line width=0.3pt] (0,0.520)--(3.200,0.520); \node[left,font=\tiny,text=muted,inner sep=1.5pt] at (0,0.520) {$2$};
\draw[muted!25,line width=0.3pt] (0,1.040)--(3.200,1.040); \node[left,font=\tiny,text=muted,inner sep=1.5pt] at (0,1.040) {$4$};
\draw[muted!25,line width=0.3pt] (0,1.560)--(3.200,1.560); \node[left,font=\tiny,text=muted,inner sep=1.5pt] at (0,1.560) {$6$};
\draw[muted!25,line width=0.3pt] (0,2.080)--(3.200,2.080); \node[left,font=\tiny,text=muted,inner sep=1.5pt] at (0,2.080) {$8$};
\draw[muted!25,line width=0.3pt] (0,2.600)--(3.200,2.600); \node[left,font=\tiny,text=muted,inner sep=1.5pt] at (0,2.600) {$10$};
\draw[-latex,muted,line width=0.6pt] (0,0)--(3.500,0) node[right,font=\scriptsize,text=navy]{X};
\draw[-latex,muted,line width=0.6pt] (0,0)--(0,2.900) node[above,font=\scriptsize,text=navy]{Y};
\fill[teal] (0.533,0.780) circle (0.085);
\fill[teal] (1.067,1.820) circle (0.085);
\fill[teal] (1.600,0.260) circle (0.085);
\fill[teal] (2.133,1.560) circle (0.085);
\fill[teal] (2.667,0.780) circle (0.085);
\end{tikzpicture}\end{center}

\ex{3} \par\smallskip\begin{center}\begin{tikzpicture}[baseline=(current bounding box.center),line cap=round]
\draw[muted!25,line width=0.3pt] (0.000,0)--(0.000,3.600); \node[below,font=\tiny,text=muted,inner sep=1.5pt] at (0.000,0) {$0$};
\draw[muted!25,line width=0.3pt] (0.500,0)--(0.500,3.600); \node[below,font=\tiny,text=muted,inner sep=1.5pt] at (0.500,0) {$1$};
\draw[muted!25,line width=0.3pt] (1.000,0)--(1.000,3.600); \node[below,font=\tiny,text=muted,inner sep=1.5pt] at (1.000,0) {$2$};
\draw[muted!25,line width=0.3pt] (1.500,0)--(1.500,3.600); \node[below,font=\tiny,text=muted,inner sep=1.5pt] at (1.500,0) {$3$};
\draw[muted!25,line width=0.3pt] (2.000,0)--(2.000,3.600); \node[below,font=\tiny,text=muted,inner sep=1.5pt] at (2.000,0) {$4$};
\draw[muted!25,line width=0.3pt] (2.500,0)--(2.500,3.600); \node[below,font=\tiny,text=muted,inner sep=1.5pt] at (2.500,0) {$5$};
\draw[muted!25,line width=0.3pt] (3.000,0)--(3.000,3.600); \node[below,font=\tiny,text=muted,inner sep=1.5pt] at (3.000,0) {$6$};
\draw[muted!25,line width=0.3pt] (3.500,0)--(3.500,3.600); \node[below,font=\tiny,text=muted,inner sep=1.5pt] at (3.500,0) {$7$};
\draw[muted!25,line width=0.3pt] (4.000,0)--(4.000,3.600); \node[below,font=\tiny,text=muted,inner sep=1.5pt] at (4.000,0) {$8$};
\draw[muted!25,line width=0.3pt] (4.500,0)--(4.500,3.600); \node[below,font=\tiny,text=muted,inner sep=1.5pt] at (4.500,0) {$9$};
\draw[muted!25,line width=0.3pt] (5.000,0)--(5.000,3.600); \node[below,font=\tiny,text=muted,inner sep=1.5pt] at (5.000,0) {$10$};
\draw[muted!25,line width=0.3pt] (0,0.000)--(5.000,0.000); \node[left,font=\tiny,text=muted,inner sep=1.5pt] at (0,0.000) {$0$};
\draw[muted!25,line width=0.3pt] (0,0.360)--(5.000,0.360); \node[left,font=\tiny,text=muted,inner sep=1.5pt] at (0,0.360) {$1$};
\draw[muted!25,line width=0.3pt] (0,0.720)--(5.000,0.720); \node[left,font=\tiny,text=muted,inner sep=1.5pt] at (0,0.720) {$2$};
\draw[muted!25,line width=0.3pt] (0,1.080)--(5.000,1.080); \node[left,font=\tiny,text=muted,inner sep=1.5pt] at (0,1.080) {$3$};
\draw[muted!25,line width=0.3pt] (0,1.440)--(5.000,1.440); \node[left,font=\tiny,text=muted,inner sep=1.5pt] at (0,1.440) {$4$};
\draw[muted!25,line width=0.3pt] (0,1.800)--(5.000,1.800); \node[left,font=\tiny,text=muted,inner sep=1.5pt] at (0,1.800) {$5$};
\draw[muted!25,line width=0.3pt] (0,2.160)--(5.000,2.160); \node[left,font=\tiny,text=muted,inner sep=1.5pt] at (0,2.160) {$6$};
\draw[muted!25,line width=0.3pt] (0,2.520)--(5.000,2.520); \node[left,font=\tiny,text=muted,inner sep=1.5pt] at (0,2.520) {$7$};
\draw[muted!25,line width=0.3pt] (0,2.880)--(5.000,2.880); \node[left,font=\tiny,text=muted,inner sep=1.5pt] at (0,2.880) {$8$};
\draw[muted!25,line width=0.3pt] (0,3.240)--(5.000,3.240); \node[left,font=\tiny,text=muted,inner sep=1.5pt] at (0,3.240) {$9$};
\draw[muted!25,line width=0.3pt] (0,3.600)--(5.000,3.600); \node[left,font=\tiny,text=muted,inner sep=1.5pt] at (0,3.600) {$10$};
\draw[-latex,muted,line width=0.6pt] (0,0)--(5.300,0) node[right,font=\scriptsize,text=navy]{mat.};
\draw[-latex,muted,line width=0.6pt] (0,0)--(0,3.900) node[above,font=\scriptsize,text=navy]{física};
\fill[teal] (2.000,1.800) circle (0.085);
\fill[teal] (2.500,1.800) circle (0.085);
\fill[teal] (3.000,2.520) circle (0.085);
\fill[teal] (3.500,2.160) circle (0.085);
\fill[teal] (4.000,2.880) circle (0.085);
\fill[teal] (4.500,3.240) circle (0.085);
\fill[teal] (1.500,1.440) circle (0.085);
\fill[teal] (4.500,0.720) circle (0.085);
\end{tikzpicture}\end{center}El punto $\res{(9,\,2)}$ es un dato atípico. Con él la correlación es positiva pero débil ($r\approx0{,}3$); sin él, es \rest{positiva y fuerte} ($r\approx0{,}94$).


% =====================================================
\seccio{4}{Coeficiente de correlación}

\ex{1} \begin{tasks}(3)
  \task \rest{Positiva y fuerte}
  \task \rest{Negativa y débil}
  \task \rest{Nula}
  \task \rest{Negativa y fuerte}
  \task \rest{Positiva moderada}
  \task \rest{Positiva perfecta}
\end{tasks}

\ex{1} \begin{tasks}(1)
  \task $r=\frac{-6}{2{,}5\cdot3}=\frac{-6}{7{,}5}=\res{-0{,}8}$: negativa y fuerte.
  \task $\sigma_{xy}=\frac{210}{8}-24=2{,}25$; \ $r=\frac{2{,}25}{6}=\res{0{,}375}$: positiva y débil.
  \task $\sigma_{xy}=\frac{420}{10}-40=2$; \ $r=\frac{2}{3{,}3}\approx\res{0{,}61}$: positiva moderada.
\end{tasks}

\ex{2} \par\smallskip\begin{center}{\renewcommand{\arraystretch}{1.15}\small\begin{tabular}{|c|c|c|c|c|}\hline $x_i$ & $y_i$ & $x_i^2$ & $y_i^2$ & $x_i\cdot y_i$\\\hline $1$ & $1$ & $1$ & $1$ & $1$\\ $2$ & $3$ & $4$ & $9$ & $6$\\ $3$ & $2$ & $9$ & $4$ & $6$\\ $4$ & $5$ & $16$ & $25$ & $20$\\ $5$ & $4$ & $25$ & $16$ & $20$\\\hline $15$ & $15$ & $55$ & $55$ & $53$\\\hline\end{tabular}}\end{center}$\bar x=\frac{15}{5}=3$, \ $\bar y=\frac{15}{5}=3$. \ $\sigma_x=\sqrt{\frac{55}{5}-3^2}=\sqrt{2}\approx 1{,}41$, \ $\sigma_y=\sqrt{\frac{55}{5}-3^2}=\sqrt{2}\approx 1{,}41$. \ $\sigma_{xy}=\frac{53}{5}-3\cdot 3=1{,}6$. \ $r=\frac{1{,}6}{1{,}41\cdot 1{,}41}\approx\res{0{,}8}$. Correlación \rest{positiva y fuerte}.

\ex{2} \par\smallskip\begin{center}{\renewcommand{\arraystretch}{1.15}\small\begin{tabular}{|c|c|c|c|c|}\hline $x_i$ & $y_i$ & $x_i^2$ & $y_i^2$ & $x_i\cdot y_i$\\\hline $2$ & $9$ & $4$ & $81$ & $18$\\ $4$ & $7$ & $16$ & $49$ & $28$\\ $6$ & $6$ & $36$ & $36$ & $36$\\ $8$ & $4$ & $64$ & $16$ & $32$\\ $10$ & $4$ & $100$ & $16$ & $40$\\\hline $30$ & $30$ & $220$ & $198$ & $154$\\\hline\end{tabular}}\end{center}$\bar x=\frac{30}{5}=6$, \ $\bar y=\frac{30}{5}=6$. \ $\sigma_x=\sqrt{\frac{220}{5}-6^2}=\sqrt{8}\approx 2{,}83$, \ $\sigma_y=\sqrt{\frac{198}{5}-6^2}=\sqrt{3{,}6}\approx 1{,}9$. \ $\sigma_{xy}=\frac{154}{5}-6\cdot 6=-5{,}2$. \ $r=\frac{-5{,}2}{2{,}83\cdot 1{,}9}\approx\res{-0{,}97}$. Correlación \rest{negativa y fuerte}.

\ex{2} \par\smallskip\begin{center}{\renewcommand{\arraystretch}{1.15}\small\begin{tabular}{|c|c|c|c|c|}\hline km & min & $x_i^2$ & $y_i^2$ & $x_i\cdot y_i$\\\hline $1$ & $10$ & $1$ & $100$ & $10$\\ $2$ & $12$ & $4$ & $144$ & $24$\\ $2{,}5$ & $15$ & $6{,}25$ & $225$ & $37{,}5$\\ $3$ & $14$ & $9$ & $196$ & $42$\\ $4$ & $20$ & $16$ & $400$ & $80$\\ $5$ & $22$ & $25$ & $484$ & $110$\\ $6$ & $25$ & $36$ & $625$ & $150$\\ $8$ & $30$ & $64$ & $900$ & $240$\\\hline $31{,}5$ & $148$ & $161{,}25$ & $3074$ & $693{,}5$\\\hline\end{tabular}}\end{center}$\bar x=\frac{31{,}5}{8}\approx 3{,}94$, \ $\bar y=\frac{148}{8}=18{,}5$. \ $\sigma_x=\sqrt{\frac{161{,}25}{8}-3{,}94^2}=\sqrt{4{,}65}\approx 2{,}16$, \ $\sigma_y=\sqrt{\frac{3074}{8}-18{,}5^2}=\sqrt{42}\approx 6{,}48$. \ $\sigma_{xy}=\frac{693{,}5}{8}-3{,}94\cdot 18{,}5\approx 13{,}84$. \ $r=\frac{13{,}84375}{2{,}16\cdot 6{,}48}\approx\res{0{,}99}$. Sí: correlación \rest{positiva y muy fuerte}, aunque no es funcional.

\ex{3} Si están en una recta que baja, la correlación es negativa perfecta: $\res{r=-1}$. Si todas tienen $y=5$, $Y$ no varía ($\sigma_y=0$) y \rest{no se puede calcular} $r$: $Y$ no depende de $X$.

\ex{3} Con $x=10,20,30,40,50$: $\sigma_x=10\sqrt2$, $\sigma_{xy}=16$ (también se multiplica por 10) y $r=\frac{16}{10\sqrt2\cdot\sqrt2}=\res{0{,}8}$, \textbf{igual}. Cambiar de unidades (de cm a mm, por ejemplo) no cambia la forma de la nube y, por tanto, tampoco $r$.


% =====================================================
\seccio{5}{Recta de regresión}

\ex{1} \begin{tasks}(1)
  \task $m=\frac{6}{4}=1{,}5$; \ $y-8=1{,}5(x-5)\Rightarrow\res{y=1{,}5x+0{,}5}$
  \task $m=\frac{-4}{5}=-0{,}8$; \ $y-20=-0{,}8(x-10)\Rightarrow\res{y=-0{,}8x+28}$
\end{tasks}

\ex{1} Tiene que pasar por $(3,\,7)$: $2\cdot3+3=9$, \ $2\cdot3+1=7$, \ $-3+4=1$. Solo $\res{y=2x+1}$.

\ex{2} Con $\bar x=3$, $\bar y=3$, $\sigma_x^2=2$ y $\sigma_{xy}=1{,}6$: $m=\frac{\sigma_{xy}}{\sigma_x^2}=\frac{1{,}6}{2}=0{,}8$; \ $y-3=0{,}8\,(x-3)\ \Rightarrow\ \res{y=0{,}8x+0{,}6}$.\par\smallskip\begin{center}\begin{tikzpicture}[baseline=(current bounding box.center),line cap=round]
\draw[muted!25,line width=0.3pt] (0.000,0)--(0.000,3.600); \node[below,font=\tiny,text=muted,inner sep=1.5pt] at (0.000,0) {$0$};
\draw[muted!25,line width=0.3pt] (0.833,0)--(0.833,3.600); \node[below,font=\tiny,text=muted,inner sep=1.5pt] at (0.833,0) {$1$};
\draw[muted!25,line width=0.3pt] (1.667,0)--(1.667,3.600); \node[below,font=\tiny,text=muted,inner sep=1.5pt] at (1.667,0) {$2$};
\draw[muted!25,line width=0.3pt] (2.500,0)--(2.500,3.600); \node[below,font=\tiny,text=muted,inner sep=1.5pt] at (2.500,0) {$3$};
\draw[muted!25,line width=0.3pt] (3.333,0)--(3.333,3.600); \node[below,font=\tiny,text=muted,inner sep=1.5pt] at (3.333,0) {$4$};
\draw[muted!25,line width=0.3pt] (4.167,0)--(4.167,3.600); \node[below,font=\tiny,text=muted,inner sep=1.5pt] at (4.167,0) {$5$};
\draw[muted!25,line width=0.3pt] (5.000,0)--(5.000,3.600); \node[below,font=\tiny,text=muted,inner sep=1.5pt] at (5.000,0) {$6$};
\draw[muted!25,line width=0.3pt] (0,0.000)--(5.000,0.000); \node[left,font=\tiny,text=muted,inner sep=1.5pt] at (0,0.000) {$0$};
\draw[muted!25,line width=0.3pt] (0,0.600)--(5.000,0.600); \node[left,font=\tiny,text=muted,inner sep=1.5pt] at (0,0.600) {$1$};
\draw[muted!25,line width=0.3pt] (0,1.200)--(5.000,1.200); \node[left,font=\tiny,text=muted,inner sep=1.5pt] at (0,1.200) {$2$};
\draw[muted!25,line width=0.3pt] (0,1.800)--(5.000,1.800); \node[left,font=\tiny,text=muted,inner sep=1.5pt] at (0,1.800) {$3$};
\draw[muted!25,line width=0.3pt] (0,2.400)--(5.000,2.400); \node[left,font=\tiny,text=muted,inner sep=1.5pt] at (0,2.400) {$4$};
\draw[muted!25,line width=0.3pt] (0,3.000)--(5.000,3.000); \node[left,font=\tiny,text=muted,inner sep=1.5pt] at (0,3.000) {$5$};
\draw[muted!25,line width=0.3pt] (0,3.600)--(5.000,3.600); \node[left,font=\tiny,text=muted,inner sep=1.5pt] at (0,3.600) {$6$};
\draw[-latex,muted,line width=0.6pt] (0,0)--(5.300,0) node[right,font=\scriptsize,text=navy]{X};
\draw[-latex,muted,line width=0.6pt] (0,0)--(0,3.900) node[above,font=\scriptsize,text=navy]{Y};
\draw[violet,line width=1.2pt] (0.000,0.360)--(5.000,3.240);
\fill[teal] (0.833,0.600) circle (0.085);
\fill[teal] (1.667,1.800) circle (0.085);
\fill[teal] (2.500,1.200) circle (0.085);
\fill[teal] (3.333,3.000) circle (0.085);
\fill[teal] (4.167,2.400) circle (0.085);
\draw[navy,line width=1pt,fill=white] (2.500,1.800) circle (0.09);
\end{tikzpicture}\end{center}

\ex{2} Con $\bar x=6$, $\bar y=6$, $\sigma_x^2=8$ y $\sigma_{xy}=-5{,}2$: $m=\frac{\sigma_{xy}}{\sigma_x^2}=\frac{-5{,}2}{8}=-0{,}65$; \ $y-6=-0{,}65\,(x-6)\ \Rightarrow\ \res{y=-0{,}65x+9{,}9}$. Cuando $x$ aumenta 1 unidad, $y$ disminuye, de media, $0{,}65$ unidades.

\ex{2} \par\smallskip\begin{center}{\renewcommand{\arraystretch}{1.15}\small\begin{tabular}{|c|c|c|c|c|}\hline Precio & Unidades & $x_i^2$ & $y_i^2$ & $x_i\cdot y_i$\\\hline $1$ & $50$ & $1$ & $2500$ & $50$\\ $2$ & $44$ & $4$ & $1936$ & $88$\\ $3$ & $40$ & $9$ & $1600$ & $120$\\ $4$ & $32$ & $16$ & $1024$ & $128$\\ $5$ & $24$ & $25$ & $576$ & $120$\\\hline $15$ & $190$ & $55$ & $7636$ & $506$\\\hline\end{tabular}}\end{center}$\bar x=\frac{15}{5}=3$, \ $\bar y=\frac{190}{5}=38$. \ $\sigma_x=\sqrt{\frac{55}{5}-3^2}=\sqrt{2}\approx 1{,}41$, \ $\sigma_y=\sqrt{\frac{7636}{5}-38^2}=\sqrt{83{,}2}\approx 9{,}12$. \ $\sigma_{xy}=\frac{506}{5}-3\cdot 38=-12{,}8$. \ $r=\frac{-12{,}8}{1{,}41\cdot 9{,}12}\approx\res{-0{,}99}$. $m=\frac{\sigma_{xy}}{\sigma_x^2}=\frac{-12{,}8}{2}=-6{,}4$; \ $y-38=-6{,}4\,(x-3)\ \Rightarrow\ \res{y=-6{,}4x+57{,}2}$. Por cada euro que sube el precio, se venden unas $6{,}4$ unidades menos.

\ex{2} Una recta a ojo puede ser la que pasa por $(3{,}5;\ 4{,}5)$ y $(0;\ 1{,}5)$: $y\approx0{,}86x+1{,}5$ (discontinua). Con la fórmula: $\sigma_x^2=\frac{91}{6}-3{,}5^2\approx2{,}92$, \ $\sigma_{xy}=\frac{110}{6}-3{,}5\cdot4{,}5\approx2{,}58$; $m=\frac{2{,}58\ldots}{2{,}92\ldots}\approx0{,}89$; \ $y-4{,}5=0{,}89(x-3{,}5)\Rightarrow\res{y\approx0{,}89x+1{,}4}$. Son muy parecidas.\par\smallskip\begin{center}\begin{tikzpicture}[baseline=(current bounding box.center),line cap=round]
\draw[muted!25,line width=0.3pt] (0.000,0)--(0.000,3.600); \node[below,font=\tiny,text=muted,inner sep=1.5pt] at (0.000,0) {$0$};
\draw[muted!25,line width=0.3pt] (0.714,0)--(0.714,3.600); \node[below,font=\tiny,text=muted,inner sep=1.5pt] at (0.714,0) {$1$};
\draw[muted!25,line width=0.3pt] (1.429,0)--(1.429,3.600); \node[below,font=\tiny,text=muted,inner sep=1.5pt] at (1.429,0) {$2$};
\draw[muted!25,line width=0.3pt] (2.143,0)--(2.143,3.600); \node[below,font=\tiny,text=muted,inner sep=1.5pt] at (2.143,0) {$3$};
\draw[muted!25,line width=0.3pt] (2.857,0)--(2.857,3.600); \node[below,font=\tiny,text=muted,inner sep=1.5pt] at (2.857,0) {$4$};
\draw[muted!25,line width=0.3pt] (3.571,0)--(3.571,3.600); \node[below,font=\tiny,text=muted,inner sep=1.5pt] at (3.571,0) {$5$};
\draw[muted!25,line width=0.3pt] (4.286,0)--(4.286,3.600); \node[below,font=\tiny,text=muted,inner sep=1.5pt] at (4.286,0) {$6$};
\draw[muted!25,line width=0.3pt] (5.000,0)--(5.000,3.600); \node[below,font=\tiny,text=muted,inner sep=1.5pt] at (5.000,0) {$7$};
\draw[muted!25,line width=0.3pt] (0,0.000)--(5.000,0.000); \node[left,font=\tiny,text=muted,inner sep=1.5pt] at (0,0.000) {$0$};
\draw[muted!25,line width=0.3pt] (0,0.450)--(5.000,0.450); \node[left,font=\tiny,text=muted,inner sep=1.5pt] at (0,0.450) {$1$};
\draw[muted!25,line width=0.3pt] (0,0.900)--(5.000,0.900); \node[left,font=\tiny,text=muted,inner sep=1.5pt] at (0,0.900) {$2$};
\draw[muted!25,line width=0.3pt] (0,1.350)--(5.000,1.350); \node[left,font=\tiny,text=muted,inner sep=1.5pt] at (0,1.350) {$3$};
\draw[muted!25,line width=0.3pt] (0,1.800)--(5.000,1.800); \node[left,font=\tiny,text=muted,inner sep=1.5pt] at (0,1.800) {$4$};
\draw[muted!25,line width=0.3pt] (0,2.250)--(5.000,2.250); \node[left,font=\tiny,text=muted,inner sep=1.5pt] at (0,2.250) {$5$};
\draw[muted!25,line width=0.3pt] (0,2.700)--(5.000,2.700); \node[left,font=\tiny,text=muted,inner sep=1.5pt] at (0,2.700) {$6$};
\draw[muted!25,line width=0.3pt] (0,3.150)--(5.000,3.150); \node[left,font=\tiny,text=muted,inner sep=1.5pt] at (0,3.150) {$7$};
\draw[muted!25,line width=0.3pt] (0,3.600)--(5.000,3.600); \node[left,font=\tiny,text=muted,inner sep=1.5pt] at (0,3.600) {$8$};
\draw[-latex,muted,line width=0.6pt] (0,0)--(5.300,0) node[right,font=\scriptsize,text=navy]{X};
\draw[-latex,muted,line width=0.6pt] (0,0)--(0,3.900) node[above,font=\scriptsize,text=navy]{Y};
\draw[accent,line width=1pt,dashed] (0.000,0.675)--(5.000,3.375);
\draw[violet,line width=1.2pt] (0.000,0.630)--(5.000,3.420);
\fill[teal] (0.714,0.900) circle (0.085);
\fill[teal] (1.429,1.800) circle (0.085);
\fill[teal] (2.143,1.350) circle (0.085);
\fill[teal] (2.857,2.700) circle (0.085);
\fill[teal] (3.571,2.250) circle (0.085);
\fill[teal] (4.286,3.150) circle (0.085);
\draw[navy,line width=1pt,fill=white] (2.500,2.025) circle (0.09);
\end{tikzpicture}\end{center}

\ex{3} Pasa por $(\bar x,\,11)$: $11=-1{,}5\bar x+20\Rightarrow\bar x=\frac{9}{1{,}5}=\res{6}$. La pendiente es negativa, y $r$ tiene el mismo signo: correlación \rest{negativa}.

\ex{3} De $r=\frac{\sigma_{xy}}{\sigma_x\sigma_y}$: $\sigma_{xy}=0{,}8\cdot2\cdot5=\res{8}$. \ $m=\frac{8}{2^2}=2$; \ $y-30=2(x-10)\Rightarrow\res{y=2x+10}$.


% =====================================================
\seccio{6}{Estimaciones}

\ex{1} $\hat y(4)=2{,}5\cdot4+10=\res{20}$; \ $\hat y(6{,}2)=2{,}5\cdot6{,}2+10=\res{25{,}5}$. Las dos son fiables.

\ex{1} \begin{tasks}(1)
  \task \rest{Sí}: correlación fuerte e interpolación.
  \task \rest{No}: correlación débil.
  \task \rest{No}: extrapolación muy lejos de los datos.
  \task \rest{Sí}: $|r|=0{,}85$ y $\bar x$ está en el centro de los datos (la estimación es $\bar y$).
\end{tasks}

\ex{2} $\hat y(3{,}5)=0{,}8\cdot3{,}5+0{,}6=\res{3{,}4}$: fiable ($r=0{,}8$ y $x$ dentro de los datos). $\hat y(20)=0{,}8\cdot20+0{,}6=\res{16{,}6}$: \rest{no fiable}: 20 está muy lejos de los datos (de 1 a 5).

\ex{2} Con los parámetros del ejercicio 24: $m=\frac{\sigma_{xy}}{\sigma_x^2}=\frac{13{,}84375}{4{,}65234375}\approx 2{,}98$; \ $y-18{,}5=2{,}98\,(x-3{,}94)\ \Rightarrow\ \res{y=2{,}98x+6{,}78}$. $\hat y(7)\approx27{,}6$ min (fiable). $\hat y(50)\approx156$ min: \rest{no es fiable}: una persona que vive a 50 km no va andando o en bici, sino en autobús o en coche.

\ex{2} $40=-6{,}4x+57{,}2\Rightarrow6{,}4x=17{,}2\Rightarrow x=\frac{17{,}2}{6{,}4}\approx\res{2{,}69\ \text{€}}$. Como la correlación es muy fuerte, la estimación es buena.

\ex{3} \begin{tasks}(1)
  \task \par\smallskip\begin{center}{\renewcommand{\arraystretch}{1.15}\small\begin{tabular}{|c|c|c|c|c|}\hline $x_i$ (°C) & $y_i$ (helados) & $x_i^2$ & $y_i^2$ & $x_i\cdot y_i$\\\hline $14$ & $20$ & $196$ & $400$ & $280$\\ $18$ & $35$ & $324$ & $1225$ & $630$\\ $20$ & $40$ & $400$ & $1600$ & $800$\\ $23$ & $60$ & $529$ & $3600$ & $1380$\\ $25$ & $65$ & $625$ & $4225$ & $1625$\\ $28$ & $80$ & $784$ & $6400$ & $2240$\\ $30$ & $95$ & $900$ & $9025$ & $2850$\\ $32$ & $105$ & $1024$ & $11\,025$ & $3360$\\\hline $190$ & $500$ & $4782$ & $37\,500$ & $13\,165$\\\hline\end{tabular}}\end{center}$\bar x=\frac{190}{8}=23{,}75$, \ $\bar y=\frac{500}{8}=62{,}5$. \ $\sigma_x=\sqrt{\frac{4782}{8}-23{,}75^2}=\sqrt{33{,}69}\approx 5{,}8$, \ $\sigma_y=\sqrt{\frac{37\,500}{8}-62{,}5^2}=\sqrt{781{,}25}\approx 27{,}95$. \ $\sigma_{xy}=\frac{13\,165}{8}-23{,}75\cdot 62{,}5=161{,}25$. \ $r=\frac{161{,}25}{5{,}8\cdot 27{,}95}\approx\res{0{,}99}$. Positiva y muy fuerte.
  \task $m=\frac{\sigma_{xy}}{\sigma_x^2}=\frac{161{,}25}{33{,}6875}\approx 4{,}79$; \ $y-62{,}5=4{,}79\,(x-23{,}75)\ \Rightarrow\ \res{y=4{,}79x-51{,}18}$
  \task $\hat y(27)\approx\res{78}$ helados; \ $\hat y(40)\approx\res{140}$. La de 27 °C, que está dentro de los datos (de 14 a 32 °C); la de 40 °C es una extrapolación.\par\smallskip\begin{center}\begin{tikzpicture}[baseline=(current bounding box.center),line cap=round]
\draw[muted!25,line width=0.3pt] (0.000,0)--(0.000,3.600); \node[below,font=\tiny,text=muted,inner sep=1.5pt] at (0.000,0) {$10$};
\draw[muted!25,line width=0.3pt] (0.714,0)--(0.714,3.600); \node[below,font=\tiny,text=muted,inner sep=1.5pt] at (0.714,0) {$15$};
\draw[muted!25,line width=0.3pt] (1.429,0)--(1.429,3.600); \node[below,font=\tiny,text=muted,inner sep=1.5pt] at (1.429,0) {$20$};
\draw[muted!25,line width=0.3pt] (2.143,0)--(2.143,3.600); \node[below,font=\tiny,text=muted,inner sep=1.5pt] at (2.143,0) {$25$};
\draw[muted!25,line width=0.3pt] (2.857,0)--(2.857,3.600); \node[below,font=\tiny,text=muted,inner sep=1.5pt] at (2.857,0) {$30$};
\draw[muted!25,line width=0.3pt] (3.571,0)--(3.571,3.600); \node[below,font=\tiny,text=muted,inner sep=1.5pt] at (3.571,0) {$35$};
\draw[muted!25,line width=0.3pt] (4.286,0)--(4.286,3.600); \node[below,font=\tiny,text=muted,inner sep=1.5pt] at (4.286,0) {$40$};
\draw[muted!25,line width=0.3pt] (5.000,0)--(5.000,3.600); \node[below,font=\tiny,text=muted,inner sep=1.5pt] at (5.000,0) {$45$};
\draw[muted!25,line width=0.3pt] (0,0.000)--(5.000,0.000); \node[left,font=\tiny,text=muted,inner sep=1.5pt] at (0,0.000) {$0$};
\draw[muted!25,line width=0.3pt] (0,0.450)--(5.000,0.450); \node[left,font=\tiny,text=muted,inner sep=1.5pt] at (0,0.450) {$20$};
\draw[muted!25,line width=0.3pt] (0,0.900)--(5.000,0.900); \node[left,font=\tiny,text=muted,inner sep=1.5pt] at (0,0.900) {$40$};
\draw[muted!25,line width=0.3pt] (0,1.350)--(5.000,1.350); \node[left,font=\tiny,text=muted,inner sep=1.5pt] at (0,1.350) {$60$};
\draw[muted!25,line width=0.3pt] (0,1.800)--(5.000,1.800); \node[left,font=\tiny,text=muted,inner sep=1.5pt] at (0,1.800) {$80$};
\draw[muted!25,line width=0.3pt] (0,2.250)--(5.000,2.250); \node[left,font=\tiny,text=muted,inner sep=1.5pt] at (0,2.250) {$100$};
\draw[muted!25,line width=0.3pt] (0,2.700)--(5.000,2.700); \node[left,font=\tiny,text=muted,inner sep=1.5pt] at (0,2.700) {$120$};
\draw[muted!25,line width=0.3pt] (0,3.150)--(5.000,3.150); \node[left,font=\tiny,text=muted,inner sep=1.5pt] at (0,3.150) {$140$};
\draw[muted!25,line width=0.3pt] (0,3.600)--(5.000,3.600); \node[left,font=\tiny,text=muted,inner sep=1.5pt] at (0,3.600) {$160$};
\draw[-latex,muted,line width=0.6pt] (0,0)--(5.300,0) node[right,font=\scriptsize,text=navy]{°C};
\draw[-latex,muted,line width=0.6pt] (0,0)--(0,3.900) node[above,font=\scriptsize,text=navy]{helados};
\draw[violet,line width=1.2pt] (0.099,0.000)--(4.874,3.600);
\draw[gold,line width=0.8pt,dashed] (2.429,0)--(2.429,1.756)--(0,1.756);
\fill[teal] (0.571,0.450) circle (0.085);
\fill[teal] (1.143,0.787) circle (0.085);
\fill[teal] (1.429,0.900) circle (0.085);
\fill[teal] (1.857,1.350) circle (0.085);
\fill[teal] (2.143,1.462) circle (0.085);
\fill[teal] (2.571,1.800) circle (0.085);
\fill[teal] (2.857,2.138) circle (0.085);
\fill[teal] (3.143,2.362) circle (0.085);
\draw[gold,line width=1.2pt] (2.429,1.756) circle (0.1);
\end{tikzpicture}\end{center}
\end{tasks}

\ex{3} \begin{tasks}(1)
  \task \par\smallskip\begin{center}{\renewcommand{\arraystretch}{1.15}\small\begin{tabular}{|c|c|c|c|c|}\hline $x_i$ & $y_i$ & $x_i^2$ & $y_i^2$ & $x_i\cdot y_i$\\\hline $1$ & $19$ & $1$ & $361$ & $19$\\ $2$ & $17$ & $4$ & $289$ & $34$\\ $3$ & $15$ & $9$ & $225$ & $45$\\ $5$ & $12$ & $25$ & $144$ & $60$\\ $9$ & $7$ & $81$ & $49$ & $63$\\\hline $20$ & $70$ & $120$ & $1068$ & $221$\\\hline\end{tabular}}\end{center}$\bar x=\frac{20}{5}=4$, \ $\bar y=\frac{70}{5}=14$. \ $\sigma_x=\sqrt{\frac{120}{5}-4^2}=\sqrt{8}\approx 2{,}83$, \ $\sigma_y=\sqrt{\frac{1068}{5}-14^2}=\sqrt{17{,}6}\approx 4{,}2$. \ $\sigma_{xy}=\frac{221}{5}-4\cdot 14=-11{,}8$. \ $r=\frac{-11{,}8}{2{,}83\cdot 4{,}2}\approx\res{-0{,}99}$. $m=\frac{\sigma_{xy}}{\sigma_x^2}=\frac{-11{,}8}{8}=-1{,}475$; \ $y-14=-1{,}475\,(x-4)\ \Rightarrow\ \res{y=-1{,}475x+19{,}9}$
  \task $\hat y(6)\approx11{,}05$ miles de euros, unos $\res{11\,050}$ €. $\hat y(20)\approx-9{,}6$: \rest{un precio negativo, que no tiene sentido}: es una extrapolación demasiado lejana (los coches no pierden valor siempre al mismo ritmo).
\end{tasks}


\end{document}
